\def\Ponderation{0}
\def\SigleNumero{STT-1920}
\def\NomCours{Méthodes statistiques}
\def\DateEvaluation{}
\def\HeureEvaluation{}
\def\Session{Hiver 2026}
\def\NomEvaluation{Révision Examen 3}

\def\Coordonnatrice{}
\def\Coordonnateur{}
\def\Enseignant{Julien Miron}
\def\Enseignante{}
\def\SectionA{}
\def\SectionB{}
\def\SectionC{}
\def\SectionD{}

\def\Directives{
\begin{enumerate}
\item Ce document sert a la revision de l'examen 3.
\item Les solutions sont affichees pour faciliter l'etude.
\item Justifiez vos reponses comme a l'examen.
\end{enumerate}
}

\newcommand{\affichepoints}{}
\providecommand{\Gradescope}{}

\documentclass{Evaluation}
\usepackage{multicol,graphicx}

\printanswers

\begin{document}

\pageblanche
\begin{questions}

\question[8] QCM rapide de concepts (une seule bonne reponse par item).
\begin{parts}
\part[2] En test a 1 echantillon, si la valeur-$p$ vaut $0{,}12$ et $\alpha=0{,}05$, on...
\begin{oneparcheckboxes}
\correctchoice ne rejette pas $H_0$.
\choice rejette $H_0$.
\choice ne peut rien conclure.
\end{oneparcheckboxes}
\begin{solution}
Comme $0{,}12>0{,}05$, on ne rejette pas $H_0$.
\end{solution}

\part[2] En ANOVA a 1 facteur, l'hypothese nulle est...
\begin{oneparcheckboxes}
\correctchoice $\mu_1=\mu_2=\cdots=\mu_a$.
\choice $\sigma_1^2=\sigma_2^2=\cdots=\sigma_a^2$.
\choice $p_1=p_2=\cdots=p_a$.
\end{oneparcheckboxes}
\begin{solution}
Le test $F$ global d'ANOVA compare l'egalite des moyennes.
\end{solution}

\part[2] En regression lineaire simple, si $\beta_1<0$, alors...
\begin{oneparcheckboxes}
\correctchoice la relation lineaire moyenne est decroissante.
\choice les variables sont independantes.
\choice la variance de $Y$ est negative.
\end{oneparcheckboxes}
\begin{solution}
$\beta_1$ est la pente de la droite moyenne.
\end{solution}

\part[2] Pour un khi-carre d'independance avec un tableau $3\times 2$, le nombre de ddl vaut...
\begin{oneparcheckboxes}
\choice 1
\correctchoice 2
\choice 3
\end{oneparcheckboxes}
\begin{solution}
$\nu=(r-1)(c-1)=(3-1)(2-1)=2$.
\end{solution}
\end{parts}

\pageblanche
\question[8] Test a 1 echantillon.
On observe un echantillon normal de taille $n=16$ avec $\bar{x}=52$ et $s=8$.
On veut tester $H_0:\mu=50$ contre $H_1:\mu>50$ au seuil $5\%$.
\begin{parts}
\part[3] Calculez la statistique de test.
\begin{solution}
\[
T_0=\frac{\bar{x}-\mu_0}{s/\sqrt{n}}=\frac{52-50}{8/4}=\frac{2}{2}=1.
\]
\end{solution}

\part[3] Donnez la region critique.
\begin{solution}
On rejette $H_0$ si $T_0>t_{0{,}95;15}\approx 1{,}753$.
\end{solution}

\part[2] Conclusion.
\begin{solution}
Comme $1<1{,}753$, on ne rejette pas $H_0$ au seuil de $5\%$.
\end{solution}
\end{parts}

\pageblanche
\question[8] ANOVA a 1 facteur.
Trois traitements sont compares avec un total de 18 observations. La table partielle est:
\begin{verbatim}
Source        SS      ddl      MS      F0
Traitement   24.0      2      ???     ???
Erreur       30.0     15      ???
\end{verbatim}
\begin{parts}
\part[3] Completez $MS_{\text{trt}}$ et $MS_E$.
\begin{solution}
$MS_{\text{trt}}=24/2=12$ et $MS_E=30/15=2$.
\end{solution}

\part[3] Calculez $F_0$.
\begin{solution}
$F_0=MS_{\text{trt}}/MS_E=12/2=6$.
\end{solution}

\part[2] Au seuil de 5\%, si $F_{0{,}05;2,15}=3{,}68$, concluez.
\begin{solution}
Comme $6>3{,}68$, on rejette $H_0$: au moins une moyenne differe.
\end{solution}
\end{parts}

\pageblanche
\question[8] Regression lineaire simple.
On donne les resultats suivants: $n=10$, $\bar{x}=4$, $\bar{y}=13$, $S_{XX}=20$, $S_{XY}=30$, et $R^2=0{,}81$.
\begin{parts}
\part[3] Calculez $\hat{\beta}_1$ et $\hat{\beta}_0$.
\begin{solution}
\[
\hat{\beta}_1=\frac{S_{XY}}{S_{XX}}=\frac{30}{20}=1{,}5,
\quad
\hat{\beta}_0=\bar{y}-\hat{\beta}_1\bar{x}=13-1{,}5(4)=7.
\]
Donc $\hat{y}=7+1{,}5x$.
\end{solution}

\part[2] Donnez la prediction ponctuelle pour $x=6$.
\begin{solution}
$\hat{y}(6)=7+1{,}5(6)=16$.
\end{solution}

\part[3] Interpretez $R^2=0{,}81$.
\begin{solution}
$81\%$ de la variabilite de $Y$ est expliquee par la relation lineaire avec $X$.
\end{solution}
\end{parts}

\pageblanche
\question[8] Khi-carre d'ajustement.
Une machine produit quatre types de capsules. Le fabricant annonce les proportions
$0{,}25, 0{,}25, 0{,}25, 0{,}25$. Sur $n=80$ capsules observees, on obtient
$O=(16,24,20,20)$.
\begin{parts}
\part[2] Donnez les frequences attendues sous $H_0$.
\begin{solution}
$E_i=80(0{,}25)=20$ pour chaque categorie.
\end{solution}

\part[4] Calculez $D_0$.
\begin{solution}
\[
D_0=\sum\frac{(O_i-E_i)^2}{E_i}
=\frac{(16-20)^2}{20}+\frac{(24-20)^2}{20}+\frac{(20-20)^2}{20}+\frac{(20-20)^2}{20}
=0{,}8+0{,}8=1{,}6.
\]
\end{solution}

\part[2] Combien de ddl?
\begin{solution}
$\nu=k-1=4-1=3$.
\end{solution}
\end{parts}

\pageblanche
\question[8] Khi-carre d'independance.
On observe le tableau suivant:
\begin{center}
\begin{tabular}{c|cc|c}
 & Oui & Non & Total\\
\hline
Groupe A & 30 & 20 & 50\\
Groupe B & 10 & 40 & 50\\
\hline
Total & 40 & 60 & 100
\end{tabular}
\end{center}
\begin{parts}
\part[2] Formulez $H_0$ et $H_1$.
\begin{solution}
$H_0$: les deux variables sont independantes.
$H_1$: elles ne sont pas independantes.
\end{solution}

\part[3] Calculez les frequences attendues.
\begin{solution}
Pour A-Oui: $E=50\times40/100=20$.
Pour A-Non: $E=50\times60/100=30$.
Pour B-Oui: $20$.
Pour B-Non: $30$.
\end{solution}

\part[3] Calculez $D_0$.
\begin{solution}
\[
D_0=\frac{(30-20)^2}{20}+\frac{(20-30)^2}{30}+\frac{(10-20)^2}{20}+\frac{(40-30)^2}{30}
=5+3{,}33+5+3{,}33=16{,}66.
\]
\end{solution}
\end{parts}

\pageblanche
\question[8] Test a 1 echantillon - objectif: intervalle de confiance.
Une variable suit une loi normale. On observe $n=25$, $\bar{x}=18$ et $s=5$.
\begin{parts}
\part[4] Donnez un IC a 95\% pour $\mu$ (utilisez $t_{0{,}025;24}=2{,}064$).
\begin{solution}
\[
IC_{95\%}:\ \bar{x}\pm t_{0{,}025;24}\frac{s}{\sqrt{n}}
=18\pm 2{,}064\times\frac{5}{5}
=18\pm 2{,}064.
\]
Donc $IC_{95\%}=[15{,}936\ ;\ 20{,}064]$.
\end{solution}

\part[4] Peut-on affirmer que $\mu=20$ est plausible avec cet IC?
\begin{solution}
Oui. La valeur 20 appartient a l'intervalle $[15{,}936\ ;\ 20{,}064]$.
\end{solution}
\end{parts}

\pageblanche
\question[8] ANOVA a 1 facteur - objectif: lire le plan experimental.
On donne la table ANOVA suivante:
\begin{verbatim}
Source        ddl
Traitement     3
Erreur        20
Total         23
\end{verbatim}
\begin{parts}
\part[4] Combien de niveaux du facteur ont ete compares?
\begin{solution}
En ANOVA a 1 facteur, $ddl_{\text{trt}}=a-1$. Donc $a=3+1=4$ niveaux.
\end{solution}

\part[4] Combien d'observations totales ont ete utilisees?
\begin{solution}
En ANOVA, $ddl_{\text{total}}=N-1$. Donc $N=23+1=24$ observations.
\end{solution}
\end{parts}

\pageblanche
\question[8] Regression lineaire simple - objectif: intervalle pour la pente.
On estime le modele $Y=\beta_0+\beta_1X+\varepsilon$ et on obtient
$\hat{\beta}_1=1{,}2$, $SE(\hat{\beta}_1)=0{,}4$, $n=12$.
\begin{parts}
\part[4] Donnez un IC a 95\% pour $\beta_1$ (utilisez $t_{0{,}025;10}=2{,}228$).
\begin{solution}
\[
\hat{\beta}_1\pm t_{0{,}025;10}SE(\hat{\beta}_1)
=1{,}2\pm 2{,}228(0{,}4)
=1{,}2\pm 0{,}891.
\]
Donc $IC_{95\%}(\beta_1)=[0{,}309\ ;\ 2{,}091]$.
\end{solution}

\part[4] La relation lineaire est-elle significative au seuil de 5\%?
\begin{solution}
Oui. L'intervalle ne contient pas 0, donc $\beta_1\neq 0$ au seuil 5\%.
\end{solution}
\end{parts}

\pageblanche
\question[8] Khi-carre d'ajustement - objectif: valider les conditions et conclure vite.
On veut verifier si une loterie produit 5 symboles equiprobables.
Sur 100 tirages: $O=(18,22,21,19,20)$.
\begin{parts}
\part[3] Donnez les frequences attendues et verifiez la condition minimale.
\begin{solution}
Sous $H_0$, chaque symbole a proba $0{,}2$, donc $E_i=100\times0{,}2=20$.
La condition $E_i\ge 5$ est satisfaite pour les 5 categories.
\end{solution}

\part[3] Combien de ddl?
\begin{solution}
$\nu=k-1=5-1=4$ (aucun parametre estime).
\end{solution}

\part[2] Si la valeur-$p$ calculee est $0{,}84$, quelle est la conclusion a 5\%?
\begin{solution}
Comme $0{,}84>0{,}05$, on ne rejette pas $H_0$.
Les donnees sont compatibles avec des proportions equiprobables.
\end{solution}
\end{parts}

\pageblanche
\question[8] Khi-carre d'independance - objectif: identifier la cellule qui contribue le plus.
On observe:
\begin{center}
\begin{tabular}{c|cc|c}
 & Reussite & Echec & Total\\
\hline
Programme A & 40 & 10 & 50\\
Programme B & 20 & 30 & 50\\
\hline
Total & 60 & 40 & 100
\end{tabular}
\end{center}
\begin{parts}
\part[3] Calculez les frequences attendues.
\begin{solution}
$E_{A,R}=50\times60/100=30$, $E_{A,E}=20$, $E_{B,R}=30$, $E_{B,E}=20$.
\end{solution}

\part[3] Calculez les contributions $\frac{(O-E)^2}{E}$ de chaque cellule.
\begin{solution}
A-Reussite: $(40-30)^2/30=3{,}33$.
A-Echec: $(10-20)^2/20=5$.
B-Reussite: $(20-30)^2/30=3{,}33$.
B-Echec: $(30-20)^2/20=5$.
\end{solution}

\part[2] Quelle(s) cellule(s) contribue(nt) le plus au khi-carre?
\begin{solution}
Les cellules A-Echec et B-Echec (contribution 5 chacune) sont les plus grandes.
\end{solution}
\end{parts}

\end{questions}
\end{document}
